% \SBtitlestyle{bar}
% \SBsubtitlestyle{none}

\begin{thebibliography}{99}
\addcontentsline{toc}{section}{{\bf Tài liệu tham khảo}\rm }%
\bibitem{rade}
H. Rademacher,
{\it Higher Mathematics from an Elementary point of view},
Birkhauser, 1983.

\bibitem{aigner}
Martin Aigner  and Gunter M. Ziegler,
{\it Proofs from the book},
Springer, 1999.

\bibitem{titu} 
T. Andreescu, R. Gelca,
{\it Mathematical olympiad challenges},  Birkhauser, 2002. 

\bibitem{Larson}
Loren C. Larson,
{\it Problem-Solving through problems},
Springer-Verlag, 1983.

\bibitem{dickson}
L. E. Dickson,
{\it New first course in the theory of equations}
John Wiley \& Sons, 1946.

\bibitem{cofman}
Judita Cofman,
{\it What to solve? Problems and Suggestions for Young Mathematicians},
Clarendon Press-Oxford, 1989.

\bibitem{nhddl}
Nguyễn Hữu Điển,
{\it Phương pháp Đirichle và ứng dụng},
NXBKHKT, 1999.

\bibitem{nhdqn}
Nguyễn Hữu  Điển,
{\it Phương pháp Quy nạp toán học},
NXBGD, 2000.

\bibitem{nhdsp}
Nguyễn Hữu  Điển,
{\it Phương pháp Số phức với hình học phẳng},
NXB ĐHQG, 2000.

\bibitem{nhdpc}
Nguyễn Hữu  Điển,
{\it Những phương pháp điển hình trong giải toán phổ thông},
NXBGD, 2001.

\bibitem{nhdct}
Nguyễn Hữu  Điển,
{\it Những phương pháp giải bài toán cực trị trong hình học},
NXBKHKT, 2001.

\bibitem{nhdst}
Nguyễn Hữu  Điển,
{\it Sáng tạo trong giải toán phổ thông},
NXBGD, 2002.

\bibitem{nhddt}
Nguyễn Hữu  Điển,
{\it Đa thức và ứng dụng},
NXBGD, 2003.

\bibitem{nhdptvd}
Nguyễn Hữu  Điển,
{\it Giải phương trình vô định nghiệm nguyên},
NXBĐHQG, 2004.

\bibitem{nhdinvar}
Nguyễn Hữu  Điển,
{\it Giải toán bằng phương pháp đại lượng bất biến},
NXBGD, 2004.


\end{thebibliography}


